% Derived from independently maintained chapters/03.tex, not the legacy fallback.
% Physical definitions and critical labels preserved; hierarchy reorganized.
\section{数据处理与公共物理模型}
\label{sec:foundation}

各问首先都需要判断一项运输任务是否能飞，以及执行它需要多少时间和能量。本章将节点、货箱、装备和地形数据转换为统一任务属性。后续改变的是任务组合、资源时序与组织方式，而不是在不同问题中任意改用另一套距离或能耗公式。

\subsection{输入数据整理与一致性检查}
\subsubsection{节点、货箱与装备数据关联}
位置表以节点编号关联，逐箱清单以货箱编号唯一标识。将同一服务区、同一物资类型的箱数汇总后与需求汇总表核对，避免重复计入汇总需求和明细需求。整理结果为1个基地O01、15个服务区和80箱物资，总质量758\,kg、总体积2.011\,m$^3$。记服务区集合为$\mathcal S$，空间节点集合为$\mathcal V=\{O01\}\cup\mathcal S$，货箱集合为$\mathcal B$。

节点$i$保留原表经纬度$(\lambda_i,\varphi_i)$与地面高程$z_i$；货箱$b$记录目的地$s(b)$、质量$\mu_b$和体积$\nu_b$。配送时限和优先权重在问题二以后进入调度评价，不参与问题一的静态组批。图\ref{fig:scene}展示位置与需求，点面积表示质量，不表示灾情严重程度。
\fig{scene.pdf}{原始地形与服务区物资需求分布。底图仅为显示重投影，模型计算仍使用原始DEM；需求点面积与质量成正比。}{fig:scene}

机型参数见表\ref{tab:common_drone}。含电池空载质量、额定载荷和电池可用能量分别参与不同公式，不能混为有效载重。准备、装载、交接、升降速度和爬升效率同样按机型读取。
\begin{table}[htbp]\centering\small
\caption{三类运输无人机主要物理参数}\label{tab:common_drone}
\begin{tabular}{lrrr}\toprule 参数 & A型 & B型 & C型\\\midrule
含电池空载质量/kg &70&65&69.9\\
额定载荷/kg &25&30&80\\
货舱体积/m$^3$ &0.060&0.073&0.250\\
巡航速度/(m/s)&12&15&15\\
空载航程/m&25000&28000&26000\\
满载航程/m&20000&16000&12000\\
可用电池能量/kWh&4.5&4.0&8.0\\
基准返航安全余量&20\%&20\%&20\%\\\bottomrule
\end{tabular}
\end{table}

\subsubsection{地形数据及空间坐标转换}
TIF与MAT两种地形文件分别核对高程矩阵、像元中心坐标和仿射边界。模型直接在原始栅格上计算航段经过地形，不用平滑或重采样改变山体。节点高程仍使用位置表值，不能用邻近DEM值悄然替换。地形显示与计算使用不同副本，显示分辨率不改变物理判定。

以O01为中心建立WGS84方位等距投影（AEQD）\cite{snyder1987}。令$\boldsymbol x_i=(x_i,y_i)^{\mathsf T}$，则
\begin{equation}
\boldsymbol x_i=\operatorname{AEQD}_{O01}(\lambda_i,\varphi_i),\qquad
 d_{ij}=\|\boldsymbol x_j-\boldsymbol x_i\|_2.
\label{eq:common_projection}
\end{equation}
任意访问节点间的水平航迹规定为投影平面直线：
\begin{equation}
\boldsymbol r_{ij}(u)=(1-u)\boldsymbol x_i+u\boldsymbol x_j,\qquad 0\le u\le1.
\label{eq:common_planar_path}
\end{equation}
以O01为端点的径向航段可用椭球测地线独立核验；这不意味着任意服务区间的投影距离都等于测地线距离。求与地形栅格的相交时，需要连续反投影式\eqref{eq:common_planar_path}，不能对经纬度端点直接线性插值代替。

\subsection{地形约束下的航段几何}
\subsubsection{航迹与地形像元的完整相交}
原始GeoTIFF采用像元中心定位。设左上参考像元中心为$(\lambda_c,\varphi_c)$，角度间隔为$\Delta\lambda,\Delta\varphi>0$，则栅格左上边界为
\begin{equation}
\lambda_0=\lambda_c-\tfrac12\Delta\lambda,\qquad
\varphi_0=\varphi_c+\tfrac12\Delta\varphi.
\label{eq:common_half_pixel}
\end{equation}
反投影航迹$(\lambda(u),\varphi(u))$对应连续列、行坐标
\begin{equation}
\xi(u)=\frac{\lambda(u)-\lambda_0}{\Delta\lambda},\qquad
\eta(u)=\frac{\varphi_0-\varphi(u)}{\Delta\varphi}.
\label{eq:common_raster_coordinates}
\end{equation}
半像元校正用于对应TIF中心定位与MAT边界变换，不是人为移动航迹。求出轨迹与整数行列边界的交点并按$u$排序；相邻交点之间没有新的栅格边界，检查其中点及全部交点即可得到相交像元集合$\mathcal C_{ij}$。采用闭像元约定，格边与角点接触的所有像元均计入。

问题一的径向航段在核实坐标单调性后逐边界求根，这一处理不外推到任意非单调长轨迹。超出DEM范围或经过无效高程时拒绝该几何输入，不以零高程填补。固定步长采样可以用于显示与对照，但不能保证发现短暂接触的每个像元。

\subsubsection{巡航海拔、作业高度与飞行阶段}
设$Z_c$为像元高程，普通运输的航段最高地形、巡航海拔及作业高度为
\begin{equation}
 z_{ij}^{\max}=\max_{c\in\mathcal C_{ij}}Z_c,\qquad H_{ij}=z_{ij}^{\max}+50,
\label{eq:common_cruise_altitude}
\end{equation}
\begin{equation}
 h_i=\begin{cases}z_i,&i=O01,\\z_i+30,&i\in\mathcal S,\end{cases}
 \quad h_{ij}^{\uparrow}=H_{ij}-h_i,\quad h_{ij}^{\downarrow}=H_{ij}-h_j.
\label{eq:common_work_height}
\end{equation}
每段均经历爬升、巡航和下降。完成服务区交接后，下一段从该区作业高度重新起算，空载返航也可能再次爬升。计算前检查$H_{ij}\ge\max(h_i,h_j)$；若不成立，应报告高度口径冲突，不能把负爬升直接截为零。中继悬停端点可能高于该巡航面，第\ref{sec:q3}章单独定义中继转场规则，不反向修改普通运输模型。

\begin{figure}[htbp]\centering
\begin{tikzpicture}[x=1cm,y=1cm,>=Stealth,font=\small]
\fill[black!5] (0,0)--(0,.25)--(1,.4)--(1.8,1.1)--(2.5,.65)--(3.4,1.5)--(4.2,.8)--(5.5,.35)--(7,.55)--(7,0)--cycle;
\draw[black!55] (0,.25)--(1,.4)--(1.8,1.1)--(2.5,.65)--(3.4,1.5)--(4.2,.8)--(5.5,.35)--(7,.55);
\draw[dashed,black!40] (0,2.25)--(7.2,2.25);
\draw[->,thick,cA] (.2,.6)--(.2,2.25);\draw[->,thick,cA] (.2,2.25)--(6.8,2.25);\draw[->,thick,cA](6.8,2.25)--(6.8,.9);
\draw[<->,cR] (3.4,1.5)--node[right]{50\,m}(3.4,2.25);
\node[left,align=right] at (.2,1.5){爬升\\$H_{ij}-h_i$};
\node[right,align=left] at (6.8,1.5){下降\\$H_{ij}-h_j$};
\node[above] at (2.2,2.25){水平巡航：$d_{ij}$};\node[above]at(6.1,2.25){$H_{ij}$};
\node[below] at (.2,0){节点$i$};\node[below]at(6.8,0){节点$j$};
\node at (3.6,-.65){各航段分别确定巡航面，本站交接后重新起算下一段};
\end{tikzpicture}
\caption{普通运输航段的三阶段飞行与高度关系（机制示意，非真实地形剖面）}\label{fig:common_leg_geometry}
\end{figure}

\subsection{运输任务的载荷与时间计算}
\subsubsection{任务组成与逐站卸货}
一项任务表示为$r=(m_r,\mathcal B_r,\pi_r)$，其中机型$m_r\in\{A,B,C\}$，箱组为$\mathcal B_r$，有序路线$\pi_r=(v_0,v_1,\ldots,v_K,v_{K+1})$满足$v_0=v_{K+1}=O01$。中间节点与箱组目的地一致，当前任务构造不重复访问同一服务区。第$j$站交付数量与质量为
\begin{equation}
 n_{rj}=|\{b\in\mathcal B_r:s(b)=v_j\}|,\qquad
 D_{rj}=\sum_{\substack{b\in\mathcal B_r\\s(b)=v_j}}\mu_b.
\label{eq:common_drop}
\end{equation}
起飞容量与逐站载荷满足
\begin{equation}
 q_{r,0}=\sum_{b\in\mathcal B_r}\mu_b\le Q_{m_r},\qquad
 \sum_{b\in\mathcal B_r}\nu_b\le V_{m_r},
\label{eq:common_capacity}
\end{equation}
\begin{equation}
 q_{r,j}=q_{r,j-1}-D_{rj},\qquad q_{r,0}\ge\cdots\ge q_{r,K}=0.
\label{eq:common_load}
\end{equation}
$v_j\to v_{j+1}$使用$q_{r,j}$计算，包括$j=K$的空载返航。途中只卸货、不补货，故初始容量合格后质量与体积不再增加；但能耗必须按卸货后的各段载荷计算。改变访问顺序会改变高载荷飞行距离，不能只比较总路程。

\subsubsection{交接完成时刻与返航时间}
设爬升、巡航、下降速度为$v_m^\uparrow,v_m^c,v_m^\downarrow$，则
\begin{equation}
 t_{ij}^{\mathrm{fly}}(m)=\frac{h_{ij}^{\uparrow}}{v_m^{\uparrow}}+\frac{d_{ij}}{v_m^c}+\frac{h_{ij}^{\downarrow}}{v_m^{\downarrow}}.
\label{eq:common_leg_time}
\end{equation}
以开始准备为任务相对零时刻。固定准备、单箱装载、每站基础交接和单箱交接时间分别为$t_m^{\mathrm{prep}},t_m^{\mathrm{load}},t_m^{\mathrm{hand}},t_m^{\mathrm{box}}$。令$n_r=|\mathcal B_r|$，起飞时刻$\theta_{r,0}=t_{m_r}^{\mathrm{prep}}+n_rt_{m_r}^{\mathrm{load}}$。到达与交接完成时间依次为
\begin{equation}
\begin{aligned}
 a_{r,j}&=\theta_{r,j-1}+t_{v_{j-1}v_j}^{\mathrm{fly}}(m_r),\\
 \theta_{r,j}&=a_{r,j}+t_{m_r}^{\mathrm{hand}}+n_{rj}t_{m_r}^{\mathrm{box}},\quad j=1,\ldots,K.
\end{aligned}
\label{eq:common_delivery_recursion}
\end{equation}
货箱以\emph{交接完成}而非抵达作为送达时刻；同站货箱共享该交接完成时刻。若$b$在第$j$站交付，则$\delta_{rb}=\theta_{r,j}$，绝对开始时刻为$s_r$时有
\begin{equation}
 C_b=s_r+\delta_{rb},\qquad p_r=\theta_{r,K}+t_{v_KO01}^{\mathrm{fly}}(m_r).
\label{eq:common_delivery_return}
\end{equation}
展开得到
\begin{equation}
\begin{split}
 p_r={}&t_{m_r}^{\mathrm{prep}}+n_rt_{m_r}^{\mathrm{load}}
 +\sum_{j=0}^Kt_{v_jv_{j+1}}^{\mathrm{fly}}(m_r)\\
 &+\sum_{j=1}^K\bigl(t_{m_r}^{\mathrm{hand}}+n_{rj}t_{m_r}^{\mathrm{box}}\bigr).
\end{split}
\label{eq:common_task_duration}
\end{equation}
该递推同时输出每箱交付偏移和返航时间，不需要按总时长平均分摊。问题一用$T^{\mathrm{sum}}=\sum_rp_r$衡量累计工作量；并行调度用$T_{\max}=\max_r(s_r+p_r)$衡量最后返航时刻。二者不能替代。

\subsection{能耗计算与返航安全约束}
\subsubsection{载荷相关航程与分阶段能耗}
配送能耗需要考虑载荷与任务组织\cite{zhang2021energy,zhang2023corrigendum}。具体计算沿用题面给定关系，不拟合其他型号的功率曲线。对$0\le q\le Q_m$，等效航程为
\begin{equation}
 L_m(q)=L_m^0-(L_m^0-L_m^F)\left(\frac q{Q_m}\right)^{3/2}.
\label{eq:common_range}
\end{equation}
该量是水平能耗的标定长度，不是一次往返可到达的单程半径。水平与爬升能耗分别为
\begin{equation}
 E_m^{\mathrm{hor}}(d,q)=E_m\frac d{L_m(q)},\qquad
 E_m^{\mathrm{up}}(q,h)=\frac{(m_m^0+q)gh}{\eta_m\,3.6\times10^6}.
\label{eq:common_energy_parts}
\end{equation}
$m_m^0$包含电池质量，$\eta_m$为爬升效率，$g=9.81\,\mathrm{m/s^2}$；焦耳除以$3.6\times10^6$转换为kWh。当前参数化中，下降计时但不另加能耗；普通交接不额外添加未给定的悬停功率。中继长时间悬停由第\ref{sec:q3}章单独计算。这是模型约定，不是对实际下降或交接能量的普遍断言。

\subsubsection{任务总能耗与返航剩余电量}
按实际逐段载荷求和：
\begin{equation}
 E_r=\sum_{j=0}^{K}\left[
 E_{m_r}\frac{d_{v_jv_{j+1}}}{L_{m_r}(q_{r,j})}
 +\frac{(m_{m_r}^0+q_{r,j})g\,h_{v_jv_{j+1}}^{\uparrow}}{\eta_{m_r}\,3.6\times10^6}\right].
\label{eq:common_task_energy}
\end{equation}
最后一段$q_{r,K}=0$，但水平飞行和空载爬升项仍然存在。每次任务使用满电电池时，返航安全判据为
\begin{equation}
 E_r\le(1-\alpha_{m_r})E_{m_r},\qquad
 SOC_r^{\mathrm{return}}=1-\frac{E_r}{E_{m_r}}\ge\alpha_{m_r}.
\label{eq:common_reserve}
\end{equation}
基准$\alpha_m=0.20$。式\eqref{eq:common_capacity}与\eqref{eq:common_reserve}共同决定物理可行性。它们本身不保证有限资源排程或连续通信可行。

\subsection{能源恢复与资源占用}
\subsubsection{两阶段充电恢复时间}
设完全充电等效时间为$\tau^{\mathrm{full}}$，任务结束荷电状态为$\rho$。从0充至90\%占完全充电时间65\%，90\%至100\%占35\%，故
\begin{equation}
 \tau(\rho)=\tau^{\mathrm{full}}\begin{cases}
 0.65(0.9-\rho)/0.9+0.35,&0\le\rho<0.9,\\
 0.35(1-\rho)/0.1,&0.9\le\rho\le1.
 \end{cases}
\label{eq:common_recharge}
\end{equation}
在$\rho=0,0.9,1$处分别得到$\tau^{\mathrm{full}}$、$0.35\tau^{\mathrm{full}}$、0，且90\%处连续。该函数不能简化为“耗电比例乘完全充电时间”。充电敏感性试验统一缩放$\tau^{\mathrm{full}}$，仍保留上述两段比例。

\subsubsection{飞机、电池与中继资源的释放区别}
令$c_r=\tau(SOC_r^{\mathrm{return}})$，从任务开始准备起，运输飞机与电池分别占用
\begin{equation}
 I_r^U=[s_r,s_r+p_r),\qquad I_r^B=[s_r,s_r+p_r+c_r).
\label{eq:common_occupancy}
\end{equation}
运输机返航后释放，原电池须充满才能再用；末次充电不计入最后返航时间，但仍计入能源占用。中继实体还有规定的返航周转，其释放不能直接套用$I_r^U$。各区间使用原始精度；同刻释放与开始允许衔接，不能先舍入分钟再计算资源峰值。

\subsection{公共模型核验}
\subsubsection{几何与单任务计算核对}
先检查$L_m(0)=L_m^0$、$L_m(Q_m)=L_m^F$、完整交付后$q_{r,K}=0$，再逐项核对时间、交付偏移、SOC及充电分段连续性。以S006的C型单点任务为例，6箱共59\,kg、0.156\,m$^3$，能耗2.597987865\,kWh、作业时长1561.553868\,s。代回8\,kWh电池容量得到返航SOC为67.525152\%，恢复时长为$0.512318349\tau_C^{\mathrm{full}}$。本例只检查任务属性，未把单架次可行误写成系统排程可行。

已有几何核验对15条基地径向航段使用独立测地线和像元条带求交：距离最大差为$1.412\times10^{-9}$\,m，相交集合一致；0.25\,m固定采样对照合计漏过4个相交像元。这说明完整相交优于该采样对照，但不构成任意服务区间投影误差的上界，也不能识别DEM分辨率以下的障碍物。本轮进一步从原始任务属性重算问题二23个架次及80箱交付，与保存计划在$10^{-6}$\,s和$10^{-7}$\,kWh范围内对应。

\subsubsection{任务输出与后续模型接口}
公共计算将任务转换为机型、箱组、路线、载荷序列、$p_r,E_r,SOC_r,c_r$和$\delta_{rb}$。问题一在此基础上选择可行批型，问题二加入有限飞机、电池及配送时限，问题三使用分段三维轨迹检查通信，问题四冻结运输与中继的实际占用区间。

对于敏感性计算，输入扰动必须进入相关物理项，而非直接放大结果工期。例如能耗增加会同时改变SOC与充电时间；交接延迟改变后续交付和返航，但在当前无额外交接功率的模型中不自动增加能耗。本文分别保留名义计算、固定决策重放和允许调整后的结果，避免把不同实验权限混为一个“通过率”。
